\subsection{A resolvable group without an expansive sequence: Problem 21.107}
\label{cand:21.107}
\problemstatement{21.107}

As usual, the open sets in the definition of an expansive
sequence are understood to be nonempty: every such open
set must meet all sufficiently late terms of a sequence
of pairwise disjoint finite subsets.
Take a free ultrafilter \(\mathcal U\) on \(\mathbb N\),
obtained by extending the cofinite filter, and the countable
Boolean group \(B=[\mathbb N]^{<\omega}\) under symmetric
difference. For \(A\in\mathcal U\), put
\(H_A=[A]^{<\omega}\) and use these subgroups as an
identity neighborhood base. Intersections are
\(H_A\cap H_C=H_{A\cap C}\), so this defines a group
topology. It is Hausdorff by deleting a coordinate
on which two elements differ, and nondiscrete because
every \(A\in\mathcal U\) is infinite.

The sets
\[
 D_k=\{s\in B:v_2(|s|+1)=k\},\qquad k\geq0,
\]
partition \(B\) into dense subsets. Given \(x+H_A\)
and \(k\), choose an arbitrarily large
\(r=2^k(2q+1)-1\geq|x|\), and add to \(x\) a
subset of \(A\setminus x\) of size \(r-|x|\).
This gives a point of the coset in \(D_k\).

Suppose \((F_n)\) is any sequence of pairwise disjoint
finite subsets of \(B\). Infinitely many empty terms
already preclude expansiveness. Otherwise discard
the finitely many empty terms and the at most one
term containing the identity, leaving an infinite
sequence of nonempty finite families of nonempty
supports. Put \(S_n=\{\max s:s\in F_n\}\).
For every \(r\), only finitely many \(S_n\) meet
\(\{1,\ldots,r\}\), since there are just \(2^r-1\)
nonempty supports with maximum at most \(r\) and
the \(F_n\) are disjoint.
Choose a subsequence with
\(\max S_{n_j}<\min S_{n_{j+1}}\).
Let \(E\) be the union of its even-indexed \(S_{n_j}\).
Either \(E\in\mathcal U\) or its complement is in
\(\mathcal U\). Whichever belongs to \(\mathcal U\)
misses all \(S_{n_j}\) of the opposite parity.
The corresponding open subgroup \(H_A\) therefore
misses infinitely many \(F_{n_j}\), a contradiction.

Thus countable resolvability does not imply existence
of an expansive sequence. Only an ordinary free
ultrafilter is used, with no selectivity assumption.
The underlying filter topology is standard; the
dense partition and obstruction are given explicitly.
Source: \artifact{research/21.107-proof.md}.

